\documentclass[11pt]{article}
\usepackage{acl2015}
\usepackage{times}
\usepackage{url}
\usepackage{latexsym}
\usepackage{natbib}
\usepackage{color}
\definecolor{citecolor}{RGB}{34,139,34}
\usepackage[pagebackref=true,breaklinks=true,letterpaper=true,citecolor=citecolor,colorlinks,bookmarks=false]{hyperref}
\usepackage[capitalise]{cleveref}
\usepackage[labelfont={bf,small},font=small]{caption}
\newcommand{\minisection}[1]{\noindent{\textbf{#1.}}}

\title{Simulating Multi-Stage Interconnection Networks in Python}

\author{Ferola \\
  Sapienza University of Rome \\
  Advanced Architectures \\
  {\tt email@domain} \\}

\date{}

\begin{document}
\maketitle

\begin{abstract}
This project builds and tests four interconnection networks used
in parallel computers: Baseline, Beneš, Clos, and XGFT. Each
network is built in Python, given a routing method, and checked
for correctness. The Clos network is also tested for blocking, and
all four are compared by size and cost.
\end{abstract}

\section{Task description}
An interconnection network connects many processors (or switches)
together inside a parallel computer. This project studies four
kinds of such networks and shows, in code, how data moves through
each one from an input to an output.

\subsection{Examples}
For a Baseline network, sending data from input 2 to output 6
follows a path such as \texttt{s0\_1 -> s1\_3 -> s2\_2}, one switch
per stage. For a Clos network, a call from input 5 to output 8
might use the path \texttt{in\_1 -> mid\_3 -> out\_2}, picking one
middle switch out of several possible ones.

\subsection{Real-world applications}
These networks are used to connect processors, memory, and servers
in supercomputers and data centers, where many devices need fast,
reliable paths to each other at once.

\section{Related work}
The Clos network was introduced by \citet{clos1953} as a way to
build large switching networks from small switches, with a proven
condition for avoiding blocking. The Beneš network, described by
\citet{benes1965}, showed that the same idea could route any full
permutation with no clashes. Fat-tree networks, including the
generalized form (XGFT) built here, are a common choice in today's
data center networks.

\section{Tools and methods}
The project is written in Python, using the \texttt{networkx}
library to build and check each network as a graph, and
\texttt{matplotlib} to draw pictures of them. No external dataset
is needed: each network is generated directly by code, with a few
example sizes used throughout (for example, Baseline and Beneš
with 8 inputs; Clos with 4 input switches, 3 inputs each, and 5
middle switches; XGFT with height 2).

\section{Evaluation approach}
In the literature, these networks are usually judged by: how many
switches and wires they need for a given size, how they are
routed, and whether they can get blocked (unable to find a free
path) under load. This project follows the same three measures.

\section{Experiments}
For each network, every possible input-output path was checked
against the actual wires built, to confirm the routing code is
correct. For Beneš, this was done using 60 random full
permutations. For Clos, correctness was also checked for the
non-blocking condition directly: with $m \geq 2n-1$ middle
switches, no random traffic pattern caused any blocked call, across
many trials.

Switch and wire counts were also measured as network size grows,
and, for Clos, as the number of middle switches changes.

\subsection{Results}
Beneš needs close to twice as many switches as Baseline for the
same size, matching the theory (Beneš has $2\log_2 N - 1$ stages,
Baseline has $\log_2 N$). Both grow with $N \log N$, not $N^2$,
which is why these networks scale well.

For Clos, blocking rate drops sharply as the number of middle
switches increases, reaching zero once $m \geq 2n - 1$, exactly as
predicted by the non-blocking theorem. Below that point, blocking
still happens sometimes, but less often than the worst-case
guarantee would suggest, since random traffic is unlikely to hit
the exact worst case.

\subsection{Discussion}
The routing methods used here are simple: Baseline and XGFT use
direct rules based on the destination address, Clos always tries
middle switches in a fixed order, and Beneš uses the Looping
Algorithm only for full permutations. None of these adapt to
traffic patterns, so a busier or more clever routing strategy could
likely do better. Blocking was only measured for Clos; the other
three networks were not tested this way.

\section{Interactive demo}
All four networks, and their routing, can also be seen in a
single self-contained web page (\texttt{network\_demo.html}). It
animates a chosen path for Baseline, Beneš, Clos, and XGFT, shows
the Looping Algorithm routing a full Beneš permutation at once, and
plays a live side-by-side comparison of Clos with too few versus
enough middle switches, so blocking can be watched happening. The
page needs no setup: results from every check are shown directly
on it.

\section{Conclusions}
This project built four common interconnection networks in Python,
routed data through each one, and checked the results against
known theory. All four networks and their routing methods worked
correctly, and the results (switch counts, blocking rates) matched
what the theory predicts. Future work could add blocking analysis
for the other networks, smarter routing strategies, or connect
these networks to an actual workload, such as an arithmetic unit,
to see how the choice of network affects real performance.

\minisection{Contributions} Single-author project: Ferola did all
the design, implementation, and testing.

\bibliographystyle{plain}
\bibliography{refs}
\end{document}
